Conversion-rate models must be trained before their labels exist. Most published corrections for
this delay assume the training stream can be revisited --- a record is written at click time and
repaired later by duplicating the event, relabelling it, or adding heads at further horizons. An
append-only pipeline cannot do that cheaply. We impose a \emph{single-visit} constraint --- at
most one training record per click, its label frozen when written --- and measure how many
published methods satisfy it: reproducing \nScreened{} of them from scratch on the Criteo
conversion logs at a one-day attribution window, \nFail{} write more than one record for at
least one click, up to \maxDup{} on average and \maxVisits{} at worst. What survives is a narrow
design space in which the only remaining choice is the target the single record carries. We
study four points on that target-hardness axis across five served capacities
(\capMin{}--\capMax{} parameters) and three pipeline write delays. Waiting for the matured label
beats training on the censored one by a wide margin. Replacing part of the hard label with an
offline teacher's estimate improves on waiting in all \cdWaitN{} cells; removing the hard label
entirely does not, crossing from better than waiting to worse as capacity grows. The hybrid
target matches an online label corrector to within noise rather
than beating it, so the case for an upstream teacher is architectural: it is what lets $K$
objectives with $K$ different attribution windows share one record written at one time.
Freezing the teacher into a build-time artifact costs \frzMean{} $\NE$, and nothing at all at
write delays of 30 minutes or more.
